\subsection{Resultados para QFT}

QFT mostró el comportamiento más sensible a la estructura de la entrada. Cuando el estado
inicial presenta periodicidad, tanto Blosc como ZFP reducen memoria de forma apreciable. En
cambio, cuando la entrada posee alta entropía y fases aleatorias, la compresión deja de ser
ventajosa y puede empeorar simultáneamente la memoria y el tiempo de ejecución. Esta
diferencia convierte a QFT en el caso más ilustrativo para discutir el alcance real del
esquema propuesto.

\subsubsection{Superposición periódica}

\begin{table}[H]
    \centering
    \small
    \setlength{\tabcolsep}{3pt}
    \caption{QFT con superposición periódica}
    \label{tab:qft-pattern-results}
    \resizebox{\linewidth}{!}{%
    \begin{tabular}{lcccccc}
        \hline
        \textbf{Qubits} & \textbf{Estrategia} & \textbf{Mem. pico (MB)} & \textbf{Ahorro mem. (\%)} & \textbf{Tiempo (s)} & \textbf{Tiempo rel. (x)} & \textbf{Error} \\
        \hline
        22 & Referencia & 73.6 & 0.0 & 0.44 & 1.00 & 0 \\
        22 & Blosc & 35.7 & 51.4 & 1.52 & 3.48 & 0 \\
        22 & ZFP & 19.0 & 74.2 & 2.52 & 5.75 & $3.79\times 10^{-11}$ \\
        23 & Referencia & 139.6 & 0.0 & 1.21 & 1.00 & 0 \\
        23 & Blosc & 34.5 & 75.3 & 3.17 & 2.61 & 0 \\
        23 & ZFP & 25.4 & 81.8 & 5.68 & 4.68 & $3.79\times 10^{-11}$ \\
        24 & Referencia & 271.6 & 0.0 & 2.75 & 1.00 & 0 \\
        24 & Blosc & 92.4 & 66.0 & 6.99 & 2.54 & 0 \\
        24 & ZFP & 45.6 & 83.2 & 12.10 & 4.41 & $3.79\times 10^{-11}$ \\
        \hline
    \end{tabular}
    }
\end{table}

En esta familia de entrada, ambas estrategias comprimidas logran ahorros espaciales
apreciables. Blosc ahorra entre 51.4\% y 75.3\% de memoria y mantiene un costo temporal entre
$2.54\times$ y $3.48\times$. ZFP amplía ese ahorro hasta el rango 74.2\%--83.2\%, pero eleva
el tiempo a entre $4.41\times$ y $5.75\times$ la referencia e introduce un error máximo
absoluto constante del orden de $10^{-11}$. En consecuencia, la compresión es útil en QFT
cuando la entrada conserva periodicidad, aunque la elección entre Blosc y ZFP depende de si se
privilegia el mayor ahorro espacial o una combinación más conservadora de tiempo y exactitud.

\subsubsection{Superposición de alta entropía}

\begin{table}[H]
    \centering
    \small
    \setlength{\tabcolsep}{3pt}
    \caption{QFT con superposición de alta entropía}
    \label{tab:qft-random-results}
    \resizebox{\linewidth}{!}{%
    \begin{tabular}{lcccccc}
        \hline
        \textbf{Qubits} & \textbf{Estrategia} & \textbf{Mem. pico (MB)} & \textbf{Ahorro mem. (\%)} & \textbf{Tiempo (s)} & \textbf{Tiempo rel. (x)} & \textbf{Error} \\
        \hline
        22 & Referencia & 71.6 & 0.0 & 0.77 & 1.00 & 0 \\
        22 & Blosc & 85.3 & -19.0 & 9.67 & 12.52 & 0 \\
        22 & ZFP & 73.1 & -2.1 & 79.52 & 103.01 & $3.53\times 10^{-11}$ \\
        23 & Referencia & 135.7 & 0.0 & 2.00 & 1.00 & 0 \\
        23 & Blosc & 152.2 & -12.1 & 20.98 & 10.48 & 0 \\
        23 & ZFP & 187.6 & -38.2 & 171.47 & 85.65 & $2.82\times 10^{-11}$ \\
        24 & Referencia & 263.6 & 0.0 & 4.33 & 1.00 & 0 \\
        24 & Blosc & 392.8 & -49.0 & 44.30 & 10.24 & 0 \\
        24 & ZFP & 409.8 & -55.5 & 376.40 & 87.03 & $3.79\times 10^{-11}$ \\
        \hline
    \end{tabular}
    }
\end{table}

La situación se invierte por completo cuando la entrada presenta alta entropía. En este caso,
ni Blosc ni ZFP ofrecen ventaja frente al almacenamiento no comprimido. El ahorro de memoria
se vuelve negativo en todos los tamaños: Blosc introduce un sobrecosto espacial entre 12.1\%
y 49.0\%, mientras que ZFP lo lleva hasta 55.5\%. El deterioro temporal es todavía más
marcado, con factores entre $10.24\times$ y $12.52\times$ para Blosc y entre $85.65\times$ y
$103.01\times$ para ZFP. Los errores de ZFP permanecen en un rango pequeño, del orden de
$10^{-11}$, de modo que el problema no es la exactitud admitida sino la falta de compresibilidad
del estado. Desde el punto de vista experimental, este es el resultado más importante de QFT:
cuando el estado carece de regularidad local, la compresión del vector completo deja de ser una
estrategia útil incluso si se tolera una pequeña discrepancia numérica.
